% Generated draft: requires mathematical and source review.
\documentclass[11pt]{article}
\usepackage[margin=25mm]{geometry}
\usepackage{fontspec}
\setmainfont{Latin Modern Roman}
\usepackage{amsmath,amssymb,mathtools,unicode-math}
\setmathfont{Latin Modern Math}
\usepackage{enumitem,needspace,xurl,hyperref}
\setlength{\parindent}{0pt}
\setlength{\parskip}{5pt}
\setlength{\emergencystretch}{3em}
\newcommand{\N}{\mathbb{N}}
\newcommand{\Z}{\mathbb{Z}}
\newcommand{\Q}{\mathbb{Q}}
\newcommand{\R}{\mathbb{R}}
\newcommand{\C}{\mathbb{C}}
\newcommand{\F}{\mathbb{F}}
\newcommand{\A}{\mathbb{A}}
\newcommand{\PP}{\mathbb{P}}
\newcommand{\E}{\mathbb{E}}
\newcommand{\eps}{\varepsilon}
\newcommand{\dd}{\,\mathrm{d}}
\begin{document}
\section*{UnsolvedMath ID: 30003040}
\textbf{Source problem code:} catalogue entry
\subsection{Definitions and mathematical statement}
Let \(k\) be a perfect field, meaning that every algebraic extension of \(k\) is separable. Choose a separable closure \(k_s\) and write \(\Gamma_k=\operatorname{Gal}(k_s/k)\). Under the usual absolute Galois cohomological dimension convention, the condition \(\operatorname{cd}(k)\leq 2\) means that \(H^n(\Gamma_k,M)=0\) for every discrete torsion \(\Gamma_k\)-module \(M\) and every integer \(n>2\). The supplied record does not specify a cohomological dimension convention in positive characteristic; this interpretation requires source confirmation.

Let \(G\) be a semisimple simply connected algebraic group over \(k\) of absolute type \(E_7\), so that its root system over an algebraic closure of \(k\) is of type \(E_7\). No splitness or isotropy assumption is supplied. The pointed set \(H^1(k,G)\) is the nonabelian continuous Galois cohomology set: its elements are continuous maps \(a\colon\Gamma_k\to G(k_s)\) satisfying
\[
 a_{\sigma\tau}=a_\sigma\,\sigma(a_\tau),
\]
modulo the equivalence \(a'_{\sigma}=g^{-1}a_\sigma\sigma(g)\) for \(g\in G(k_s)\). Its distinguished element is the class of the constant identity cocycle.

The supplied type-\(E_7\) formulation of Serre's conjecture II asks whether
\[
 \text{for every perfect field }k\text{ with }\operatorname{cd}(k)\leq 2,
 \quad\text{for every such }G/k,
 \qquad H^1(k,G)=\{1\}.
\]
Equivalently, must every \(G\)-torsor over \(k\) have a \(k\)-rational point?

The record labels this problem open. Its supplied literature assessment discusses remaining exceptional cases, including \(E_8\), but does not specifically establish the status of this type-\(E_7\) formulation. That assessment and the intended cohomological dimension convention require independent mathematical review.
\subsection{Short English statement}
For every perfect field of cohomological dimension at most two and every semisimple simply connected group of type \(E_7\) over that field, must every torsor under the group be trivial? The supplied open label is not independently verified, and the accompanying general literature assessment does not settle the status of this specific formulation.
\subsection{Sources}
\begin{itemize}
\item Supplied local UnsolvedMath registry record, permanent ID \(30003040\), source problem code \texttt{catalogue entry}, in the Oberwolfach Reports Open Problems collection. The record attributes the conjecture to Serre through its title; no individual proposer of this workshop problem is named. The supplied proposed year is 2016.
\item \emph{Algebraic Cobordism and Projective Homogeneous Varieties} (2016), Oberwolfach Reports, pp.~215--265, DOI 10.4171/owr/2016/5. Supplied source report and literature reference labelled ``original OWR report'': \url{https://doi.org/10.4171/owr/2016/5}. No problem-specific page or passage is supplied.
\item Supplied literature reference labelled ``Gille survey on Serre Conjecture II'': \url{https://math.univ-lyon1.fr/homes-www/gille/publis/hyderabad.pdf}.
\item Literature assessment dated exactly \texttt{2026-08-21}: supplied dated curation, not independently verified. It states that Serre Conjecture II is proved in several classical and exceptional cases and that unresolved exceptional cases include \(E_8\) and related remaining cases, and consequently labels the extracted general claim open. The supplied background explicitly calls this dated literature triage derived from OWR curation and requests independent verification of time-sensitive status and cited claims. This assessment does not specifically substantiate the open status of the type-\(E_7\) statement above.
\item The registry's supplied statement status is \texttt{corrected\_verified}; its verification note says only that notation, typography, and sentence structure were repaired without changing mathematical content. This editorial label is not independent mathematical or source verification.
\end{itemize}
\end{document}
